% GENERATED FILE. Edit canonical lessons, then run npm run generate:books.
% canonical-source-hash: fnv1a64:7f9368960933824c
% canonical-lessons: FR-C334-point-de-vue, FR-C334-soutenir, FR-C334-pretendre, FR-C334-affirmer, FR-C334-convaincant

\chapter{Mon point de vue}
\label{ch:fr-mon-point-de-vue}
\hlchaptermodality{334}

\begin{chapteropening}
\textbf{By the end of this chapter:} \emph{I can give my point of view, maintain or claim something, state it firmly and say whether an answer is convincing.}

Complete the last lesson of chapter 334: I can give my point of view, maintain or claim something, state it firmly and say whether an answer is convincing.
\end{chapteropening}

\section[le point de vue]{le point de vue — the point of view}

\label{lesson:FR-C334-point-de-vue}



\begin{quote}
Before the new one: say the French for a team, then the French for a birthday.
\end{quote}

\subsection*{You'll want to know: le point de vue}
\textbf{le point de vue} — \textquotedblleft{}the point of view\textquotedblright{}.

\begin{quote}
De mon point de vue, c'est une bonne idée.
\end{quote}

\begin{quote}
\emph{From my point of view, it is a good idea.}
\end{quote}

The plural is \emph{les points de vue}. \emph{À mon avis} says much the same thing.

\begin{cousinweb}
\emph{point} is inherited from Latin \emph{punctum} \textquotedblleft{}a prick, a point\textquotedblright{} (EtymDB; Lewis \& Short), and \emph{vue} is the noun made from \emph{voir} \textquotedblleft{}to see\textquotedblright{}, inherited from Latin \emph{vidēre}: the spot you look from.
\end{cousinweb}

\begin{grammarlens}[title={What you've built}]
One more word for telling, explaining and describing at B1.
\end{grammarlens}

\subsection*{Your turn}
\begin{itemize}
\raggedright
  \item \emph{Say it:} \emph{le point de vue}
  \item \emph{Say it:} \emph{le point de vue}, once more
  \item \emph{Say it:} a sentence of your own with \emph{le point de vue}
  \item \emph{Recall:} say the French for a team, then the French for a birthday, then say \emph{le point de vue} again
\end{itemize}

\begin{tcolorbox}[breakable,colback=teal!4,colframe=teal!35!black,title={Before you move on}]
What does le point de vue mean? (\textquotedblleft{}The point of view\textquotedblright{}.) Say it once more.
\end{tcolorbox}
\section[soutenir]{soutenir — to maintain, to claim; to support}

\label{lesson:FR-C334-soutenir}



\begin{quote}
Before the new one: say the French for a birthday, then the French for the point of view.
\end{quote}

\subsection*{You'll want to know: soutenir}
\textbf{soutenir} — \textquotedblleft{}to maintain, to claim; to support\textquotedblright{}.

\begin{quote}
Il soutient que le train est plus confortable que l'avion.
\end{quote}

\begin{quote}
\emph{He maintains that the train is more comfortable than the plane.}
\end{quote}

Conjugated like \emph{tenir}: \emph{je soutiens, nous soutenons}; past participle \emph{soutenu}. It also means to support a person: \emph{soutenir un ami}.

\begin{cousinweb}
\emph{soutenir} is inherited from Latin \emph{sustinēre} \textquotedblleft{}to hold up, to bear\textquotedblright{} (Lewis \& Short), through Old French \emph{sustenir} (EtymDB). The Latin is built from \emph{sub} \textquotedblleft{}under\textquotedblright{} and \emph{tenēre} \textquotedblleft{}to hold\textquotedblright{}: you hold an idea up so it does not fall.
\end{cousinweb}

\begin{grammarlens}[title={What you've built}]
One more word for telling, explaining and describing at B1.
\end{grammarlens}

\subsection*{Your turn}
\begin{itemize}
\raggedright
  \item \emph{Say it:} \emph{soutenir}
  \item \emph{Say it:} \emph{soutenir}, once more
  \item \emph{Say it:} a sentence of your own with \emph{soutenir}
  \item \emph{Recall:} say the French for a birthday, then the French for the point of view, then say \emph{soutenir} again
\end{itemize}

\begin{tcolorbox}[breakable,colback=teal!4,colframe=teal!35!black,title={Before you move on}]
What does soutenir mean? (\textquotedblleft{}To maintain, to claim; to support\textquotedblright{}.) Say it once more.
\end{tcolorbox}
\section[prétendre]{prétendre — to claim (with a hint of doubt)}

\label{lesson:FR-C334-pretendre}



\begin{quote}
Before the new one: say the French for the point of view, then the French for to maintain.
\end{quote}

\subsection*{You'll want to know: prétendre}
\textbf{prétendre} — \textquotedblleft{}to claim (with a hint of doubt)\textquotedblright{}.

\begin{quote}
Il prétend qu'il n'a pas reçu mon message.
\end{quote}

\begin{quote}
\emph{He claims he did not get my message.}
\end{quote}

A false friend: \emph{prétendre} is to claim, and the speaker often doubts the claim. Conjugated like \emph{rendre}: \emph{il prétend}.

\begin{cousinweb}
\emph{prétendre} comes from Latin \emph{praetendere} \textquotedblleft{}to stretch forth, to hold out in front\textquotedblright{} (Lewis \& Short), taken into Old French as a learned word (EtymDB). You hold a claim out in front of you.
\end{cousinweb}

\begin{grammarlens}[title={What you've built}]
One more word for telling, explaining and describing at B1.
\end{grammarlens}

\subsection*{Your turn}
\begin{itemize}
\raggedright
  \item \emph{Say it:} \emph{prétendre}
  \item \emph{Say it:} \emph{prétendre}, once more
  \item \emph{Say it:} a sentence of your own with \emph{prétendre}
  \item \emph{Recall:} say the French for the point of view, then the French for to maintain, then say \emph{prétendre} again
\end{itemize}

\begin{tcolorbox}[breakable,colback=teal!4,colframe=teal!35!black,title={Before you move on}]
What does prétendre mean? (\textquotedblleft{}To claim (with a hint of doubt)\textquotedblright{}.) Say it once more.
\end{tcolorbox}
\section[affirmer]{affirmer — to state firmly, to assert}

\label{lesson:FR-C334-affirmer}



\begin{quote}
Before the new one: say the French for to maintain, then the French for to claim.
\end{quote}

\subsection*{You'll want to know: affirmer}
\textbf{affirmer} — \textquotedblleft{}to state firmly, to assert\textquotedblright{}.

\begin{quote}
Elle affirme qu'elle était à la maison ce soir-là.
\end{quote}

\begin{quote}
\emph{She states that she was at home that evening.}
\end{quote}

Stronger than \emph{dire}: you say it as a fact.

\begin{cousinweb}
\emph{affirmer} is a learned borrowing from Latin \emph{affirmāre} \textquotedblleft{}to assert, to declare\textquotedblright{} (EtymDB; Lewis \& Short), built on \emph{firmus} \textquotedblleft{}firm\textquotedblright{}: you make your words firm.
\end{cousinweb}

\begin{grammarlens}[title={What you've built}]
One more word for telling, explaining and describing at B1.
\end{grammarlens}

\subsection*{Your turn}
\begin{itemize}
\raggedright
  \item \emph{Say it:} \emph{affirmer}
  \item \emph{Say it:} \emph{affirmer}, once more
  \item \emph{Say it:} a sentence of your own with \emph{affirmer}
  \item \emph{Recall:} say the French for to maintain, then the French for to claim, then say \emph{affirmer} again
\end{itemize}

\begin{tcolorbox}[breakable,colback=teal!4,colframe=teal!35!black,title={Before you move on}]
What does affirmer mean? (\textquotedblleft{}To state firmly, to assert\textquotedblright{}.) Say it once more.
\end{tcolorbox}
\section[convaincant]{convaincant — convincing}

\label{lesson:FR-C334-convaincant}



\begin{quote}
Before the new one: say the French for to claim, then the French for to state firmly.
\end{quote}

\subsection*{You'll want to know: convaincant}
\textbf{convaincant} — \textquotedblleft{}convincing\textquotedblright{}.

\begin{quote}
Sa réponse n'était pas très convaincante.
\end{quote}

\begin{quote}
\emph{His answer was not very convincing.}
\end{quote}

It agrees: \emph{un ton convaincant, une réponse convaincante}. It comes from \emph{convaincre}, to convince.

\begin{cousinweb}
\emph{convaincant} is the present participle of \emph{convaincre}, which comes from Latin \emph{convincere} \textquotedblleft{}to overcome, to prove wrong\textquotedblright{} (Lewis \& Short), built on \emph{vincere} \textquotedblleft{}to conquer\textquotedblright{}.
\end{cousinweb}

\begin{grammarlens}[title={What you've built}]
One more word for telling, explaining and describing at B1.
\end{grammarlens}

\subsection*{Your turn}
\begin{itemize}
\raggedright
  \item \emph{Say it:} \emph{convaincant}
  \item \emph{Say it:} \emph{convaincant}, once more
  \item \emph{Say it:} a sentence of your own with \emph{convaincant}
  \item \emph{Recall:} say the French for to claim, then the French for to state firmly, then say \emph{convaincant} again
\end{itemize}

\begin{tcolorbox}[breakable,colback=teal!4,colframe=teal!35!black,title={Before you move on}]
What does convaincant mean? (\textquotedblleft{}Convincing\textquotedblright{}.) Say it once more.
\end{tcolorbox}
